\section{Scope, conventions and the question asked}
\label{sec:scope}

\subsection{The question}

The parent program asks whether a physical response can realize the shifted completed-zeta transfer
\begin{equation}
H(p)=\xi(\tfrac12+p),\qquad K_\omega(p)=\frac{H(p-\omega)}{H(p+\omega)},
\label{eq:target}
\end{equation}
as a causal, passive input/output map in an ordinary signal norm. Any candidate must first be a physical
channel with a defined input, output and energy balance; only then does a comparison make sense, and the
comparison proceeds through ordered tests---a parameter dictionary, an identity endpoint, a power front
$(2\pi/p)^\omega$ with its logarithmic tangent, the complete local response, and a hierarchy of delays at
$\log 2,\log 3,\dots$ with prescribed coefficients---each of which stops the comparison on failure. The
theory is not to be adjusted to the target. This paper reports what \Nfour{} super-Yang--Mills theory,
through its half-BPS Wilson line, offers at each stage.

The line defect was chosen because its displacement operator has a protected two-point normalization,
the Bremsstrahlung function $B(\lambda,N)$ \cite{CHMS}, so that a displacement-driven response is
computable at every coupling. The Loewner equation entered because a growing trace is a one-parameter
family of shapes that could carry a time hierarchy. Both channels were tested; neither carries the
memory the target requires. The paper records the physics found along the way.

\subsection{Conventions}

The theory is $SU(N)$ \Nfour{} SYM, $N\ge2$, $\theta=0$, with Euclidean action
$(1/g^2)\int\tr(\tfrac12F_{\mu\nu}F_{\mu\nu}+D_\mu\Phi_ID_\mu\Phi_I+\dots)$, Hermitian fields,
$\tr T^aT^b=\delta^{ab}/2$, $F_{\mu\nu}=\partial_\mu A_\nu-\partial_\nu A_\mu-\ii[A_\mu,A_\nu]$ and
$D_\mu=\partial_\mu-\ii[A_\mu,\cdot]$. Feynman-gauge propagators are $g^2\delta^{ab}\delta_{\mu\nu}/(4\pi^2x^2)$
for the gauge field and the same for each scalar. The Maldacena--Wilson line is
\begin{equation}
W=\frac1N\tr P\exp\!\int\big(\ii A\cdot\dot x+|\dot x|\,n\cdot\Phi\big),
\end{equation}
with larger parameters to the left. The straight half-BPS line is $x_0(\sigma)=(\sigma,0,0,0)$ with
constant $n_0\in S^5$. Defect correlators $\dvev{O_1\cdots O_n}$ are path-ordered insertions divided by
$\vev W$. We write $\lambda=g^2N$, $B_1=\lambda(1-1/N^2)/(16\pi^2)$ for the one-loop Bremsstrahlung
function, and use the exact inputs of \cite{CHMS}: $C_D=12B$, $C_\Phi=2B$, the cusp anomalous dimension
$\Gamma_{\rm cusp}=-B(\varphi^2-\theta^2)$ at small angles, and the radiated energy $\Delta E=2\pi B\int\dot v^2$.
At strong coupling $B\to\sqrt\lambda/4\pi^2$ and $2\pi B\to\sqrt\lambda/2\pi$, Mikhailov's coefficient
\cite{Mikhailov}.

The Loewner traces are generated by the chordal Loewner equation with driver $u$; for the linear driver
$u=at$ the prefix at time $t$ is $\sqrt t$ times the prefix generated by the slope $s=a\sqrt t$ up to time one.
The base of the trace is the origin and the tip is $q(t)$.

\subsection{Classification}

Every statement below is classified as a proof, a proof sketch, a labelled numerical computation, an
observation, a reading or a literature input, in the ledger at the end of each section. Numerical
diagnostics are floating checks of analytic statements, not interval certificates; the standard-library
programs and their records are listed in Appendix~\ref{app:validation}. Formulas taken from the
literature through automated fetch summaries were validated by limits or by symmetry before use, and
this is stated where it applies.
